\section{Related Work}
\label{sec:related}

\para{Prompt injection and indirect attacks.} Prompt injection exploits the inability of LLMs to distinguish between instructions and data~\citep{greshake2023not}. \citet{gozverev2025can} formalized this as the instruction-data separation problem and introduced the \sep benchmark, showing that current LLMs fundamentally struggle with this distinction. Defenses include data transformation techniques such as spotlighting~\citep{hines2024spotlighting}, preference optimization via SecAlign~\citep{chen2024secalign}, and architectural separation of instruction and data pathways~\citep{wallace2025aside}. Our work measures the instruction-data confusion mechanism directly using a benign proxy task, complementing these defense-oriented studies.

\para{Long-context attacks.} The \ninja framework~\citep{shah2025ninja} is the most directly relevant prior work. \citet{shah2025ninja} demonstrated that embedding harmful goals within longer benign documents increases attack success rates across Llama-3.1-8B, Qwen2.5-7B, Mistral-7B, and Gemini 2.0 Flash, with ASR rising from 23.7\% to 58.8\% for Llama as context grows from 0 to 15K words. They also showed that beginning-of-context placement and topically relevant filler text amplify attacks. Our study extends this investigation to frontier API models (GPT-4.1-mini, Gemini 2.5 Flash, Claude Sonnet 4) and reveals that the monotonically increasing trend does not generalize across all model families.

\para{Attention and long contexts.} The ``Lost in the Middle'' phenomenon~\citep{liu2024lost} showed that LLMs exhibit a U-shaped attention curve over long contexts, performing well on information at the beginning and end but poorly on information in the middle. This positional bias provides one mechanism for long-context attacks: adversarial content can exploit low-attention regions to avoid detection. Our position analysis partially corroborates this finding, with \gptmini showing a significant primacy effect (beginning placement yields 93.1\% \isr vs.\ 79.4\% for middle).

\para{Adversarial suffix optimization.} Complementary to context-based attacks, suffix optimization methods~\citep{andriushchenko2024jailbreaking} achieve near-100\% attack success on safety-aligned LLMs using short adversarial suffixes (${\sim}$25 tokens). Interestingly, \citet{andriushchenko2024jailbreaking} found a U-shaped relationship between suffix length and attack success, in contrast to the monotonically increasing relationship found by \citet{shah2025ninja} for context length. This distinction reflects different mechanisms: suffix attacks rely on adversarial gradient signals that degrade with length, while context attacks exploit attention dilution that grows with length. Our work addresses the context-based setting.

\para{Defense scaling laws.} \citet{fu2025short} discovered a square-root scaling law: adversarial training on attacks of length $k$ provides defense against attacks of length up to $O(k^2)$. This implies that the attacker's advantage from longer contexts is sublinear, and defenders need not anticipate the maximum attack length. Our finding that \claudesonnet achieves complete immunity (0\% \isr across all conditions) is consistent with this perspective---sufficiently strong safety training can neutralize the length advantage entirely.
